\section{Mechanism and discussion}
\label{sec:discussion}

% TODO(human): voice pass; consider merging with results at venue length.

\paragraph{Sparse tugs vs.\ budget allocation.} Each significant pair
back-propagates through at most a handful of sampled points, so the loss
delivers \emph{precise but sparse} pulls, iterated by re-pairing --- the toy
probes show surgery ``walking'' around a defect a few points at a time. The
resampling prior is the opposite: broad, indirect, budget-limited. The
results sort cleanly along this axis: 1-D loops, which the prior could only
over- or under-tessellate, respond to precise value pressure (H1 win); wide
2-D shells benefit from both (H2: prior helps, loss helps more, composition
best); 0-D component structure needs neither --- any spreading pressure
suffices (H0 null, three channels).

\paragraph{Why the control damages topology.} The norm-matched repulsion
imposes a length-scale floor on the sampled surface; the void's death time
pins at that equilibrium (identical across seeds) and thin structures
collapse ($\beta_2$ loss on the cube, $\beta_1$ loss on the torus). This
sharpens the control argument: spending the same gradient budget through the
same channel \emph{without} topological information is not neutral --- it is
harmful --- so C1's wins cannot be an artifact of extra regularization.

\paragraph{Measurability at the loss's density.} A persistence loss can only
enforce features that clear the significance floor at its sampling density
$M$: the torus void and the double torus's two tube loops do not, at
$M{=}2048$. We treat this as a principle rather than a nuisance: the target
bundle is built (and audited) at the same $M$, features below the floor are
excluded from the objective, and the double-torus stretch case is evaluated
against its two measurable loops --- which is also all the reconstruction
itself can express at feasible budgets. Raising $M$ buys resolution at
linear-ish persistence cost; the trade is explicit and logged at bundle-build
time.

\paragraph{Recruitment is landscape-guided in practice --- and it is the
mechanism for loops.} Its location blindness --- pull the nearest bar
anywhere --- did not misfire on our geometry because unproductive tears
self-extinguish (their bar's death saturates at the local detour scale)
while productive cuts run away toward the target. The caveat stands for
target features larger than any achievable detour; the probe and its
mechanism are reported so the term is used with eyes open. The C6 ablation
sharpens its role beyond cold-start repair: a loop that never clears the
live significance floor at $M$ receives metric pressure \emph{only} through
recruitment, at every refresh --- matching alone reduces the torus to
baseline. Recruitment is how the loss reaches below its own significance
floor on the live side.

\paragraph{Prescription.} Correct topology in the \emph{loss}; allocate
budget with a \emph{wide} prior; combine for the best of both. Curriculum
schedules are unnecessary once the loss scale is calibrated against the
photometric gradient.
